\documentclass[a4paper,10pt]{report}\input{../0.Préambule/Préambule_cours_prof}\begin{document}

\definecolor{titlecolor}{rgb}{0.0, 0.48, 0.24}
\newpage
\setcounter{chapter}{14}
\chapter{Les parallélogrammes}

\section{Aspect graphique}

\vspace{-6mm}
\noindent \begin{minipage}{11cm}
\begin{dft}
Un parallélogramme est un quadrilatère qui a ses côtés opposés parallèles deux à deux. 
\end{dft}
\medskip

\begin{rmq}
$ABCD$ est un parallélogramme, mais $ABDC$ n'en est pas un. Attention à l'ordre des lettres. 
\end{rmq}
\end{minipage}
\begin{minipage}{6cm}
\begin{center}
\definecolor{ffqqqq}{rgb}{1.,0.,0.}
\definecolor{cqcqcq}{rgb}{0.7529411764705882,0.7529411764705882,0.7529411764705882}
\begin{tikzpicture}[line cap=round,line join=round,>=triangle 45,x=0.7cm,y=0.7cm]
\clip(1.5,1.5) rectangle (9.5,7.5);
\draw [line width=1.6pt] (2.,4.)-- (7.,7.);
\draw [line width=1.6pt] (7.,7.)-- (9.,5.);
\draw [line width=1.6pt] (9.,5.)-- (4.,2.);
\draw [line width=1.6pt] (4.,2.)-- (2.,4.);
\begin{scriptsize}
\draw [color=black] (2.,4.)-- ++(-2.5pt,-2.5pt) -- ++(5.0pt,5.0pt) ++(-5.0pt,0) -- ++(5.0pt,-5.0pt);
\draw[color=black] (2.2,4.67) node {\large{$A$}};
\draw [color=black] (7.,7.)-- ++(-2.5pt,-2.5pt) -- ++(5.0pt,5.0pt) ++(-5.0pt,0) -- ++(5.0pt,-5.0pt);
\draw[color=black] (7.64,7.2) node {\large{$B$}};
\draw [color=black] (9.,5.)-- ++(-2.5pt,-2.5pt) -- ++(5.0pt,5.0pt) ++(-5.0pt,0) -- ++(5.0pt,-5.0pt);
\draw[color=black] (8.9,4.3) node {\large{$C$}};
\draw [color=black] (4.,2.)-- ++(-1.5pt,-1.5pt) -- ++(3.0pt,3.0pt) ++(-3.0pt,0) -- ++(3.0pt,-3.0pt);
\draw[color=black] (3.38,1.8) node {\large{$D$}};
\end{scriptsize}
\end{tikzpicture}
\end{center}
\end{minipage}

\vspace{-3mm}
\section{Construction d'un parallélogramme}

\vspace{-3mm}
\begin{meth}[\href{https://www.youtube.com/watch?v=7-HH2Y6oDaQ}{Avec une règle et une équerre}]

\begin{center}
\includegraphics[scale=0.48]{../40.Image/15. Les parallélogrammes 1.png}
\end{center}
\end{meth}

\begin{meth}[\href{https://www.youtube.com/watch?v=oszYYwCRaAI}{Avec une règle et un compas}]

\begin{center}
\includegraphics[scale=0.48]{../40.Image/15. Les parallélogrammes 2.png}
\end{center}
\end{meth}

\section{Caractérisation d'un parallélogramme}

\vspace{-3mm}
\noindent \begin{minipage}{11cm}
\begin{prop}
\textbf{Si} un quadrilatère est un parallélogramme, \textbf{alors} les côtés opposés ont la même longueur.
\end{prop}

\begin{prop}
\textbf{Si} un quadrilatère est un parallélogramme, \textbf{alors} ses diagonales se coupent en leur milieu.
\end{prop}
\end{minipage}
\begin{minipage}{6cm}
\begin{center}
\definecolor{ffqqqq}{rgb}{1.,0.,0.}
\definecolor{cqcqcq}{rgb}{0.7529411764705882,0.7529411764705882,0.7529411764705882}
\begin{tikzpicture}[line cap=round,line join=round,>=triangle 45,x=0.7cm,y=0.7cm]
\clip(1.5,1.5) rectangle (9.5,7.5);
\draw [domain=1.5:9.5] plot(\x,{(--14.--3.*\x)/5.});
\draw [domain=1.5:9.5] plot(\x,{(-2.--3.*\x)/5.});
\draw [domain=1.5:9.5] plot(\x,{(--28.-2.*\x)/2.});
\draw [domain=1.5:9.5] plot(\x,{(--12.-2.*\x)/2.});
\draw [line width=1.2pt,dash pattern=on 1pt off 1pt on 2pt off 4pt,color=ffqqqq] (2.,4.)-- (9.,5.);
\draw [line width=1.2pt,dash pattern=on 1pt off 1pt on 2pt off 4pt,color=ffqqqq] (7.,7.)-- (4.,2.);
\draw [line width=1.6pt] (2.,4.)-- (7.,7.);
\draw [line width=1.6pt] (7.,7.)-- (9.,5.);
\draw [line width=1.6pt] (9.,5.)-- (4.,2.);
\draw [line width=1.6pt] (4.,2.)-- (2.,4.);
\begin{scriptsize}
\draw [color=black] (2.,4.)-- ++(-2.5pt,-2.5pt) -- ++(5.0pt,5.0pt) ++(-5.0pt,0) -- ++(5.0pt,-5.0pt);
\draw[color=black] (2.2,4.67) node {\large{$A$}};
\draw [color=black] (7.,7.)-- ++(-2.5pt,-2.5pt) -- ++(5.0pt,5.0pt) ++(-5.0pt,0) -- ++(5.0pt,-5.0pt);
\draw[color=black] (7.64,6.89) node {\large{$B$}};
\draw[color=black] (-2.72,1.53) node {\large{$f$}};
\draw [color=black] (9.,5.)-- ++(-2.5pt,-2.5pt) -- ++(5.0pt,5.0pt) ++(-5.0pt,0) -- ++(5.0pt,-5.0pt);
\draw[color=black] (8.7,4.59) node {\large{$C$}};
\draw[color=black] (-2.72,-1.67) node {\large{$g$}};
\draw[color=black] (4.84,9.43) node {\large{$h$}};
\draw[color=black] (-2.72,8.53) node {\large{$i$}};
\draw [color=black] (4.,2.)-- ++(-1.5pt,-1.5pt) -- ++(3.0pt,3.0pt) ++(-3.0pt,0) -- ++(3.0pt,-3.0pt);
\draw[color=black] (3.38,2.19) node {\large{$D$}};
\draw [color=black] (5.5,4.5)-- ++(-1.5pt,-1.5pt) -- ++(3.0pt,3.0pt) ++(-3.0pt,0) -- ++(3.0pt,-3.0pt);
\draw[color=black] (5.86,4.17) node {\large{$I$}};
\end{scriptsize}
\end{tikzpicture}
\end{center}
\end{minipage}

\begin{prop}[Centre de symétrie]
\textbf{Si} un quadrilatère est un parallélogramme, \textbf{alors} il a un centre de symétrie qui est le point $I$ d'intersection de ses diagonales.
\end{prop}

\noindent \begin{minipage}{12cm}
\begin{prop}
\textbf{Si} un quadrilatère est un parallélogramme, \textbf{alors} ses angles opposés ont la même mesure.
\end{prop}

\medskip
\begin{rmq}
Les rectangles, les losanges et les carrés sont des parallélogrammes particuliers .
\end{rmq}
\end{minipage}
\begin{minipage}{5cm}
\begin{center}
\definecolor{ffqqqq}{rgb}{1.,0.,0.}
\definecolor{cqcqcq}{rgb}{0.7529411764705882,0.7529411764705882,0.7529411764705882}
\begin{tikzpicture}[line cap=round,line join=round,>=triangle 45,x=0.7cm,y=0.7cm,scale=0.9]
\clip(1.5,1.5) rectangle (9.5,7.5);
\draw [line width=1.6pt] (2.,4.)-- (7.,7.);
\draw [line width=1.6pt] (7.,7.)-- (9.,5.);
\draw [line width=1.6pt] (9.,5.)-- (4.,2.);
\draw [line width=1.6pt] (4.,2.)-- (2.,4.);
\begin{scriptsize}
\draw [color=black] (2.,4.)-- ++(-2.5pt,-2.5pt) -- ++(5.0pt,5.0pt) ++(-5.0pt,0) -- ++(5.0pt,-5.0pt);
\draw[color=black] (2.2,4.67) node {\large{$A$}};
\draw [color=black] (7.,7.)-- ++(-2.5pt,-2.5pt) -- ++(5.0pt,5.0pt) ++(-5.0pt,0) -- ++(5.0pt,-5.0pt);
\draw[color=black] (7.64,7.2) node {\large{$B$}};
\draw [color=black] (9.,5.)-- ++(-2.5pt,-2.5pt) -- ++(5.0pt,5.0pt) ++(-5.0pt,0) -- ++(5.0pt,-5.0pt);
\draw[color=black] (8.9,4.3) node {\large{$C$}};
\draw [color=black] (4.,2.)-- ++(-1.5pt,-1.5pt) -- ++(3.0pt,3.0pt) ++(-3.0pt,0) -- ++(3.0pt,-3.0pt);
\draw[color=black] (3.38,1.8) node {\large{$D$}};
\end{scriptsize}
\end{tikzpicture}
\end{center}
\end{minipage}

\vspace{-3mm}
\section{Reconnaitre un parallélogramme}

\vspace{-3mm}
\begin{meth}[La démonstration]
Il faut savoir les reconnaître, mais aussi savoir justifier pourquoi ce sont bien des parallélogrammes. Il est
important de donner les bons arguments en s'appuyant sur les informations dont on est sûr : données de
l'énoncé, codages sur le dessin ou éléments prouvés avant...

En résumé : 

\begin{center}
\boxed{\begin{minipage}{14cm}
$ABCD$ est un quadrilatère non croisé.
\begin{enumerate}
\item[•] \textbf{Si} $ABCD$ a ses côtés opposés parallèles $2$ à $2$, \textbf{alors} c'est un parallélogramme.
\item[•] \textbf{Si} $ABCD$ a ses côtés opposés égaux $2$ à $2$, \textbf{alors} c'est un parallélogramme.
\item[•] \textbf{Si} $ABCD$ a ses diagonales qui se coupent en leur milieu, \textbf{alors} c'est un parallélogramme.
\item[•] \textbf{Si} $ABCD$ a deuc côtés opposés parallèles et égaux, \textbf{alors} c'est un parallélogramme.
\end{enumerate}
\end{minipage}}
\end{center}
\end{meth}

\vspace{-3mm}
\section{Aire d'un parallélogramme}

\vspace{-3mm}
\noindent \begin{minipage}{9cm}
\begin{prop}
L'aire d'un parallélogramme est égale au produit d’un de ses côtés par la hauteur relative à ce
côté, exprimés dans la même unité :


\begin{center}
\Large{$\mathcal{A} = c \times h$}
\end{center}
$\mathcal{A}$ est l'aire du parallélogramme\\
$c$ est la mesure d'un des cotés du parallélogramme\\
$h$ est la mesure de la hauteur relative au coté $c$.
\end{prop}
\end{minipage}
\hspace{5mm}
\begin{minipage}{7cm}
\begin{center}
\begin{tikzpicture}[line cap=round,line join=round,>=triangle 45,x=1.0cm,y=1.0cm,scale=0.8]
\clip(-0.5,0.5) rectangle (8.5,5.7);
\draw [line width=1.8pt,color=blue] (2.,5.)-- (8.,5.);
\draw [line width=1.8pt,color=blue] (8.,5.)-- (6.,2.);
\draw [line width=1.8pt,color=blue] (6.,2.)-- (0.,2.);
\draw [line width=1.8pt,color=blue] (0.,2.)-- (2.,5.);
\draw [<->,line width=1.2pt] (5.,5.)-- (5.,2.);
\draw [<->,line width=1.2pt] (0.,1.5)-- (6.,1.5);
\begin{scriptsize}
\draw [color=black] (2.,5.)-- ++(-3.0pt,-3.0pt) -- ++(6.0pt,6.0pt) ++(-6.0pt,0) -- ++(6.0pt,-6.0pt);
\draw[color=black] (1.6343270187369738,5.403074949472213) node {\Large{$A$}};
\draw [color=black] (8.,5.)-- ++(-3.0pt,-3.0pt) -- ++(6.0pt,6.0pt) ++(-6.0pt,0) -- ++(6.0pt,-6.0pt);
\draw[color=black] (8.266319101774872,5.427729195134064) node {\Large{$B$}};
\draw [color=black] (6.,2.)-- ++(-3.0pt,-3.0pt) -- ++(6.0pt,6.0pt) ++(-6.0pt,0) -- ++(6.0pt,-6.0pt);
\draw[color=black] (6.343287940150501,1.7912279600110523) node {\Large{$C$}};
\draw [color=black] (0.,2.)-- ++(-3.0pt,-3.0pt) -- ++(6.0pt,6.0pt) ++(-6.0pt,0) -- ++(6.0pt,-6.0pt);
\draw[color=black] (-0.3256855113801746,1.7295923458564249) node {\Large{$D$}};
\draw[color=black] (4.568182252497234,3.566333647664319) node {\Large{$h$}};
\draw[color=black] (3.0026376529696996,1.0392734673245991) node {\Large{$c$}};
\end{scriptsize}
\end{tikzpicture}
\end{center}
\end{minipage}

\end{document}